第17.4节的分区迁移在一个已应用索引处冻结完整状态。图像里既有已经形成的原结果，也有尚未完成的C/T、会话准入、退款占额和跨分区事务责任。消费交接状态、原回执及其来源与当前履约进度也在同一图像中保存，恢复后的扫描据此判断哪些原EV已经交接。分块存储保存这一整体，业务恢复再从同一个前缀继续。只导出页面上的购买列表，会遗漏关闭旧会话、继续原付款或约束退款所依赖的事实。

\subsection{清单同时连接来源和字节}

图像清单称为manifest。它记录格式版本、逻辑分区、来源代次、应用索引、序列化格式、整幅图像的长度和摘要，以及有顺序的块列表。每个块条目记录块摘要、长度和可读副本或确切对象版本。恢复者按列表重建字节流，核对长度与摘要，再使用指定格式解码。

\CNineteenFigure{ch19-06}{manifest关联完整业务状态、来源前缀与传输字节}{fig:19-manifest}

逻辑分区b1标记这些事实的稳定归属，G1则可以是当时承载b1的副本组。来源代次说明这份图像由哪一轮合法权威形成，应用索引说明它包含哪个已应用前缀。摘要说明取得的字节与清单是否一致。保存来源回执或经过授权的发布记录以后，恢复者才能把一串校验通过的字节关联到有权产生该前缀的来源。

序列化格式也进入这份关联。同一对象的字段顺序、数字表示或转义方式改变，语义可能相同，字节摘要却会变化。若原权威对自己的结构编码计算摘要，传输层又把字段排序后重新编码，就分别保留原回执摘要与传输摘要。恢复时先核对传输字节，再按原格式或明确的语义规则核对权威回执。比较不同编码域的摘要相等，会把合法重编码误判为内容损坏。

\CNineteenFigure{ch19-24-digests}{来源回执和传输编码各有摘要，重建后再核对完整状态关系}{fig:19-digests}

完整图像里的操作映射继续以四分量
\[
S=(\texttt{account\_id},\texttt{action},\texttt{session\_id},\texttt{idempotency\_key})
\]
作为键。每条操作保存规范化原输入、P/C/T、原平台范围与\code{pay-C}、原受理事实及已经形成的原结果。会话集合、未决工作、退款、库存准备与协调事务各有保存位置。导入后的同一$S$首先找到原绑定，再决定回放原受理、返回原结果或继续未完责任。

\subsection{先保存块，再发布完整集合}

把内容摘要作为块键，可以让重复上传具有明确含义。服务重新计算收到字节的摘要，同一键已有同样内容时返回已有块，声称同一摘要却发送不同内容时拒绝。这里依赖所选摘要在当前工程假设下可靠区分不同字节，访问控制决定谁能上传与发布，加密决定哪些主体能读取内容。

每个块达到要求的保存条件后，发布者才提交manifest。例如要求两个独立副本时，发布前就确认两个副本可读且校验通过。发布服务在一次原子元数据更新中，把入口的发布代次从$g$推进为$g+1$，并安装新清单。两个发布者都从$g$开始时，其中一个先成功，另一个条件失败后重读入口。

\CNineteenFigure{ch19-07}{块上传、清单发布与旧图像回收分别有明确时点}{fig:19-publication}

发布前退出时，旧入口继续指向旧完整图像，新块暂时没有已发布引用。后台在确认没有发布责任、恢复责任或其他保留引用以后回收它们。若先让入口引用新清单，再上传最后几个块，读取者会取得一个实际缺块的图像，发布服务在安装入口前便要排除这条路径。

manifest已经提交而回复丢失时，发布者按原图像身份查询。当前入口若就是自己提交的清单，恢复原发布结果；若另一个合法版本已经推进入口，则按保留策略查找自己的历史版本。超时改变调用方掌握的信息，已提交的元数据继续确定发布事实。

发布代次和业务应用索引分别比较。发布代次仲裁并发更新，应用索引约束业务前缀的正常前进。在同一来源代次内，索引360之后的正常发布不能退到350。灾难恢复有意选择旧前缀时，由单独的恢复控制流程登记回退，暂停新身份分配，并保留与独立平台核对更晚效果的责任。

\subsection{从固定清单恢复到可运行分区}

恢复者先固定一份manifest的确切版本，核对来源及前缀，再逐块下载。某副本损坏时换同一块的有效副本。全部块的长度和摘要通过，整幅图像的长度和摘要也通过以后，才进入解码与业务结构检查。块可读、格式可解析、业务关联完整，依次提供不同层次的证据。

结构检查核对原结果里的购买编号与收款编号是否关联原受理，待办是否引用原C/T与原平台键，事件是否引用原O/EV，退款是否关联原收款并保留未知占额。关闭会话与截止信息也随绑定保留。格式升级改变字段含义时，导入器按明确转换规则处理原版本，缺失字段不能默认解释为从未受理。

\CNineteenFigure{ch19-08}{固定manifest并验证完整图像以后，目标才进入资格交接与激活}{fig:19-restore}

图像加载完成，目标组再取得第17.4节规定的承载资格。迁移依次完成旧组冻结、完整导入、目录交接和目标激活，原工作者从原C/T继续。本书的参考迁移策略在存在prepared库存或未完成协调责任时暂缓交接，由恢复者先按原决定闭合责任，再冻结前缀。

灾难恢复到旧图像时，外部平台可能已经保存更晚的付款。恢复者读取原scope和\code{pay-C}，查询原操作，汇合原结果。若恢复者新造C及平台键，就会向平台提交另一项操作。原图像里保存的稳定关联，使本地旧前缀能够与独立效果域重新对齐。

图像保留也进入发布结构。旧manifest只要承担备份、审计或在途恢复责任，它引用的块就继续保留。回收先确定可删除的清单版本，再计算已经失去全部保留引用的块，并为未完成发布留出恢复窗口。同一内容块可以被多份图像共享，删除一份清单时，引用计数或等价的可达性检查决定块能否删除。
\FloatBarrier
